\section{Numerical methods and validation}
\label{app:num-methods}
This appendix specifies the numerical calculations used in
Sec.~\ref{sec:num-examples} and in the primary-state check of
Table~\ref{tab:num-primary}. The calculations evaluate finite-chain
reduced states directly; the analytical CFT predictions enter only
the subsequent comparisons and the fixed-power fitting models.

\subsection{States, subsystems, and sampling}
\label{app:num-sampling}
We use periodic chains with unit lattice spacing and contiguous
subsystems $A=\{0,\ldots,\ell-1\}$, $B=A^c$. All active mixture
fits and figures use $\ell=6,7,8,9,10$. The Ising states are
$I$, $\sigma$, and $\varepsilon$; the boson states are
$00$, $ss$, and $3s,s$, with $s=1/\sqrt2$. Each of the three
unordered pairs is evaluated in both directions. We use natural
logarithms and the incoherent mixture
$\mu_{ab}^C=(\rho_a^C+\rho_b^C)/2$, for $C=A,B$.

The reference chain-size set is
\begin{equation}
 \mathcal N=\{32,48,64,82,96,112,128,144,160,176,192,208,224,240,256\}.
 \label{eq:num-reference-grid}
\end{equation}
The primary-to-primary check retains this reference set, giving
51 geometries per direction at $x\leq0.08$. The mixture study
uses the denser sets below, which reach $N=8192$ and $16384$.
The rounding in $\mathcal D$ is to the nearest integer; the
prefactor two makes every chain size even. Unions remove duplicate
sizes. No additional selection is made by the observed entropy.
\begin{equation}\mathcal D=\{2\operatorname{round}[30\,100^{j/48}]:j=0,\ldots,48\},\qquad \mathcal R=\{200+4j:j=0,\ldots,200\}.\label{eq:num-dense-grid}\end{equation}
\begin{align}
\mathcal N_{\rm ising}&=\mathcal N\cup\mathcal D\cup\mathcal R\cup\{4096+256j:j=0,\ldots,16\}\nonumber\\
&\quad\cup\{152,1024,1118,1146,1232,1356,1492,1528,1604,1644\}\nonumber\\
&\quad\cup\{1808,1852,1972,2048,2190,2354,2412,2654,2890,3000\}\nonumber\\
&\quad\cup\{3140,3216,3294,3436,3500,3552,3606,3660,3784,3854\}\nonumber\\
&\quad\cup\{3926,4000,4158,4222,4286,4426,4670,4734,4798,5000\}\nonumber\\
&\quad\cup\{5246,5310,5542,5758,6270,6526,6782,7038,7294,7550\}\nonumber\\
&\quad\cup\{7806,8000,8016,8032,8048,8064,8080,8096,8112,8128\}\nonumber\\
&\quad\cup\{8144\},\nonumber\\
\mathcal N_{\rm boson}&=\mathcal N\cup\mathcal D\cup\{100,104,168,292,322,354,428,472,512,762\}\nonumber\\
&\quad\cup\{924,1024,1118,1356,1492,1644,1808,2048,2412,2922\}\nonumber\\
&\quad\cup\{3540,3750,3916,4096,4234,4376,4608,4720,4834,5000\}\nonumber\\
&\quad\cup\{5222,5626,5810,6250,6466,6688,6918,7156,7402,7656\}\nonumber\\
&\quad\cup\{7920,8192,8436,8688,8948,9216,9462,9714,9974,10240\}\nonumber\\
&\quad\cup\{10488,10740,10998,11264,11764,12288,12800,13312,13814,14336\}\nonumber\\
&\quad\cup\{14840,15360,15864,16384\}.
\label{eq:num-extended-grid}\end{align}
All indices $j$ in these set definitions are integers within the stated bounds.


For each group, observations are selected individually by their
fraction $\ell/N$: a chain need not contribute all five lengths
inside the fitting window. Repeated fractions arising from
different $(N,\ell)$ remain distinct observations.
Table~\ref{tab:num-windows} gives the windows and counts.
The figures show only fitted observations. All lengths share circular
markers and colors; residual colors identify models.
Previously computed lengths
$4,5,11,12$ are excluded from these fits and figures.
\begin{table}[htbp]\centering\small
\begin{tabular}{llrrrr}\toprule
Observable & CFT & Upper cutoff & Points & Sizes & Min. training\\\midrule
$D^A$ & Ising & $0.02$ & 1255 & 278 & 1250\\
$D^A$ & Boson & $0.1$ & 598 & 124 & 593\\
$\delta$ & Ising & $0.002$ & 260 & 67 & 255\\
$\delta$ & Boson & $0.005$ & 317 & 68 & 312\\\bottomrule
\end{tabular}
\caption{Sampling per ordered pair for all three fitted models.
The lower endpoint is zero and is open. ``Min. training'' is the
smallest number of observations left after omitting a complete
chain size. These are deterministic lattice data, not independent
random samples.}
\label{tab:num-windows}
\end{table}

\subsection{Exact Ising primary states and interval matrices}
\label{app:num-ising}
We use the critical periodic spin Hamiltonian
\begin{equation}
 H=-\sum_{j=0}^{N-1}X_jX_{j+1}-\sum_{j=0}^{N-1}Z_j,
 \qquad X_N=X_0,\qquad P=\prod_j Z_j.
 \label{eq:num-ising-H}
\end{equation}
The Jordan--Wigner Majoranas
$\gamma_{2j}=(\prod_{k<j}Z_k)X_j$ and
$\gamma_{2j+1}=(\prod_{k<j}Z_k)Y_j$ reduce this Hamiltonian
to a quadratic form~\cite{numMontes2017,numPeschelEisler}.
Even spin parity gives antiperiodic fermions. Their vacuum
represents $I$, and occupying the two modes at momenta
$\pm\pi/N$ represents $\varepsilon$. Odd spin parity gives
periodic fermions; their lowest state with the zero-mode
occupation fixed to $P=-1$ represents $\sigma$. All three
states have zero spin momentum. Their exact energies are
\begin{equation}
 E_I=-2\csc\frac{\pi}{2N},\qquad
 E_\sigma=-2\cot\frac{\pi}{2N},\qquad
 E_\varepsilon=E_I+8\sin\frac{\pi}{2N}.
 \label{eq:num-ising-energies}
\end{equation}
With velocity two, these gaps approach the scaling dimensions
$0,1/8,1$, respectively.

For reproducibility, the real antisymmetric covariance is
$G_{rs}=\langle i\gamma_r\gamma_s\rangle$ for $r\ne s$ and
$G_{rr}=0$. Set $L_M=2N$, $k_m=2\pi(m+t)/L_M$, and
$w_m=-\operatorname{sgn}(\sin k_m)$. The twist is $t=1/2$
for $I,\varepsilon$ and $t=0$ for $\sigma$. Reverse $w_m$
at $m=0,N-1,N,L_M-1$ for $\varepsilon$. For $\sigma$, set
$w_0=w_N=0$ and include the parity-fixing term in
\begin{equation}
 (G_a)_{rs}=\frac{i}{L_M}\sum_{m=0}^{L_M-1}w_m e^{ik_m(r-s)}
 +\begin{cases}[(-1)^r-(-1)^s]/L_M,&a=\sigma,\\0,&a=I,\varepsilon.
 \end{cases}
 \label{eq:num-covariance}
\end{equation}
This exact Fourier construction avoids a full Hamiltonian
eigensolve. We check antisymmetry, purity $G_a^2=-\mathbb I$,
and the energies in Eq.~\eqref{eq:num-ising-energies}.

Restrict $G_a$ to the first $2\ell$ Majoranas. If an orthogonal
change of basis gives blocks
$\left(\begin{smallmatrix}0&g_j\\-g_j&0\end{smallmatrix}\right)$,
and $b_r$ are the correspondingly rotated Majoranas, then
\begin{equation}
 \rho_a^A=2^{-\ell}\prod_{j=0}^{\ell-1}
       (\mathbb I+g_j i b_{2j}b_{2j+1}).
 \label{eq:num-ising-rdm}
\end{equation}
All Jordan--Wigner strings lie inside this interval, so this is
also its spin density matrix. We form the $2^\ell$-dimensional
component matrices and then their mixture. Although each
component is Gaussian, the mixture generally is not; its
entropy is evaluated from the actual matrix, not an averaged
covariance.

\subsection{Compact-boson states and exact environment contraction}
\label{app:num-boson}
The boson uses radius $R=\sqrt2$, charges
$(Q_L,Q_R)=(a,b)/\sqrt2$, and weights
$(h,\bar h)=(a^2,b^2)/4$. The charge codes
$(a,b)=(0,0),(1,1),(3,1)$ label the three states. On sites
$z_j=e^{2\pi i j/N}$, the auxiliary vertex-block construction
with charges at zero and $W$ gives a Jastrow polynomial
\cite{numHerwerth}. For an occupied set $O$ its normalized
amplitude is
\begin{equation}
 \Psi_{a,b}(O)=\frac{1}{\sqrt{\mathcal Z_{a,b}}}
 \prod_{j\in O}z_j^{a+1}(1-z_j/W)^b
 \prod_{i<j\in O}(z_j-z_i)^2,
 \quad |O|=\frac{N-a-b}{2},\quad W=10^8.
 \label{eq:num-boson-wavefunction}
\end{equation}
$\mathcal Z_{a,b}$ is the sum of squared amplitudes before
normalization. The three states have occupations $N/2,N/2-1,N/2-2$
and are exactly orthogonal. No orthogonalization is performed.
The complex phase of this amplitude is retained in every
off-diagonal matrix element. The finite endpoint $W$ is part
of the supplied lattice-state convention; it approximates the
strict insertion at infinity.

To compute both density matrices and cross-state transitions,
let $U,V\subset A$ be occupied sets and let $Y\subset B$ be
their common environment. For states $i=1,2$, with occupations
$M_i$ and unnormalized polynomials $P_i$, we need
\begin{equation}
 (T_{12}^A)_{UV}=\frac{1}{\sqrt{\mathcal Z_1\mathcal Z_2}}
 \sum_{\substack{Y\subset B\\|Y|=M_1-|U|=M_2-|V|}}
       P_1(U\cup Y)P_2(V\cup Y)^*.
 \label{eq:num-boson-transition}
\end{equation}
It vanishes if the two required environment occupations differ.
Setting the states equal gives $\rho_a^A$.

We contract this sum through confluent Vandermonde identities
and Pfaffians, building on the moment method for the
Haldane--Shastry state~\cite{numStephan}. Explicitly, set
$m=M_1-|U|=M_2-|V|$, $\nu_j=\mathbf1_{j\in U}+\mathbf1_{j\in V}$,
$d=M_1+M_2$, $g_i(z)=(1-z/W)^{b_i}$, and
\begin{align}
 L_i(U)&=\prod_{j\in U}z_j^{a_i+1}g_i(z_j)
                  \prod_{j<k\in U}(z_j-z_k)^2,\nonumber\\
 w(z)&=z^{a_1-a_2-2M_2+2}g_1(z)\overline{g_2(z)},\nonumber\\
 \mathcal E_{UV}&=\sum_{\substack{Y\subset B\\|Y|=m}}
       \prod_{i<j\in Y}(z_i-z_j)^4
       \prod_{i\in Y}\left[w(z_i)\prod_{j\in A}(z_j-z_i)^{2\nu_j}\right],
 \label{eq:num-environment}
\end{align}
so that $(T_{12}^A)_{UV}=L_1(U)\overline{L_2(V)}
\prod_{j\in V}z_j^{-2m}\mathcal E_{UV}/\sqrt{\mathcal Z_1\mathcal Z_2}$.
Let $f(z)=(1,z,\ldots,z^{d-1})^T$ and define
\begin{equation}
 \mathsf M_B=\sum_{j\in B}w(z_j)
 [f(z_j)f'(z_j)^T-f'(z_j)f(z_j)^T].
 \label{eq:num-moment}
\end{equation}
At each site of $A$, a matrix $E_{UV}$ contains no column if
$\nu_j=0$, $f(z_j)$ if $\nu_j=1$, and $f(z_j),f'(z_j)$ if
$\nu_j=2$, in site order. With $q=\sum_j\nu_j$ and
$\Delta_{\rm ext}=\prod_{i<j\in A}(z_j-z_i)^{\nu_i\nu_j}$,
\begin{equation}
 \mathcal E_{UV}=\frac{(-1)^{q(q-1)/2}}{\Delta_{\rm ext}}
   \operatorname{Pf}\begin{pmatrix}\mathsf M_B&E_{UV}\\-E_{UV}^T&0\end{pmatrix}.
 \label{eq:num-bordered-pf}
\end{equation}
Here $d+q$ is even, including for transitions with odd $d$.
The formula defines the contraction without requiring an
inverse of a possibly singular cross-moment matrix. Production
evaluation uses sparse full-circle moments, a rank-$2\ell$
subtraction of the interval sites, and cached small Pfaffian
minors, with exact borders or regularizers for singular moments.
This extends the norm contraction to off-diagonal reduced
matrices and transitions; it is a numerical construction, not
a new CFT entropy formula.

Cancellation in the moment algebra requires 100 decimal digits
through $N=256$ and 120 digits at the added sizes. Sampled
entries are independently checked at 160 digits using direct
pivoted Pfaffians. Only the final reduced matrices are converted
to double precision. The exponentially large global vector is
never formed in production.

\subsection{Large intervals from a small Gram matrix}
\label{app:num-gram}
A global pure state has Schmidt rank at most $2^\ell$ across
the cut, even though $B$ has $N-\ell$ sites. The mixture on $B$
therefore has rank at most $2^{\ell+1}$. Its calculation needs
the transition $T_{ab}^A=\tr_B\ket a\bra b$, as well as
$\rho_a^A,\rho_b^A$; the individual reduced states alone do
not specify their Schmidt-support overlap on $B$.

Writing $\ket a=\sum_{u,v}X_a(v,u)\ket u_A\ket v_B$ and
$Z=(X_a,X_b)/\sqrt2$ gives $ZZ^\dagger=\mu_{ab}^B$ and
\begin{equation}
 \mathcal G=Z^\dagger Z=\frac12
 \begin{pmatrix}(\rho_a^A)^T&(T_{ab}^A)^*\\
 (T_{ab}^A)^T&(\rho_b^A)^T\end{pmatrix}.
 \label{eq:num-gram}
\end{equation}
The nonzero eigenvalues $\lambda_k$ of $\mathcal G$ equal
those of $\mu_{ab}^B$. For its normalized eigenvectors $v_k$
and the first-block projector $P_a$, the directed entropy is
\begin{equation}
 D^B_{a\mid b}=\tr(\rho_a^A\log\rho_a^A)
 -2\sum_{\lambda_k>0}\lambda_k(v_k^\dagger P_av_k)\log\lambda_k.
 \label{eq:num-gram-entropy}
\end{equation}
Purity supplies the first term, and the second follows from
$X_a=\sqrt2ZE_a$, where $E_a$ embeds the first block.
The reversed direction uses $P_b$ and $\rho_b^A$.
We diagonalize common parity sectors on $B$ for Ising and
common occupation sectors for the boson. The full Gram
dimension is at most $2048$ for the retained lengths.
A constant phase on $T_{ab}^A$ is immaterial, but its
configuration-dependent phases must be retained.

For Ising, direct Gaussian transition formulas between the
orthogonal primaries would divide by zero overlap. We instead
choose a local Hermitian unitary $U_A=\gamma_0=X_0$ for pairs
containing $\sigma$, or $U_A=i\gamma_0\gamma_1=-Z_0$ for
$I,\varepsilon$, and set $\ket{b'}=U_A\ket b$.
This gives a nonzero anchor overlap on the sampled grid.
Gaussian product identities~\cite{numBravyi,numSurace} give
\begin{align}
 K&=\mathbb I-G_aG_{b'},&
 o&=2^{-N/2}|\det K|^{1/4},\nonumber\\
 KZ_A&=(\mathbb I-iG_a)E_A,&
 G_{*,A}&=i\left[E_A^T(\mathbb I-iG_{b'})Z_A-\mathbb I_{2\ell}\right],
 \label{eq:num-anchor-solve}
\end{align}
where $E_A$ selects the first $2\ell$ coordinates. Fourier
application of $G_a$ and a pivoted real LU factorization of
$K$ avoid forming a full complex inverse; its log determinant
also gives $\log o$. For a complex antisymmetric $H$ on the
interval modes, Wick reconstruction is
\begin{equation}
 \mathcal R(H)=2^{-\ell}
 \sum_{\substack{S\subset\{0,\ldots,2\ell-1\}\\|S|\ {
m even}}}
 i^{|S|/2}\operatorname{Pf}(H_S)
       \prod_{r\in S}^{\longrightarrow}\gamma_r.
 \label{eq:num-wick}
\end{equation}
The product is ordered by increasing index and the empty
Pfaffian is one. Restoring the anchor yields
$T_{ab}^A=o\mathcal R(G_{*,A})U_A$, up to a single irrelevant
phase. Independent small-chain wavefunction checks validate
this use of the Gaussian identities. The boson transitions
are obtained directly from Eq.~\eqref{eq:num-boson-transition}.

\subsection{Entropy evaluation and control of cancellation}
\label{app:num-precision}
For ordinary fractions, diagonalizing the component and mixture
matrices gives $S(\rho\|\mu)=\tr\rho\log\rho-\tr\rho\log\mu$.
The cross trace uses the weights of $\rho$ in the eigenbasis
of $\mu$; the matrices need not commute. The primary check
uses the same spectral trace with $\mu$ replaced by $\rho_b^A$
and a target-logarithm floor of $10^{-15}$.

At very small fractions, subtracting nearly equal entropy traces,
or subtracting $D^B$ from $\log2$, loses relative precision.
We evaluate the remainder directly using matrix-function
divided differences~\cite{numHigham}. Let $f(t)=t\log t$,
$D=\operatorname{diag}(d_i)>0$, and
$H=D+E=V\operatorname{diag}(\lambda_k)V^\dagger$. Then
\begin{align}
 \mathcal R_i&=[f(H)-f(D)]_{ii}-f'(d_i)E_{ii}\nonumber\\
 &=\sum_k|(V^\dagger E)_{ki}|^2
 \frac{\lambda_k\log(\lambda_k/d_i)-\lambda_k+d_i}
      {(\lambda_k-d_i)^2}.
 \label{eq:num-stable-remainder}
\end{align}
The identity follows by expanding in the eigenvectors of $H$
and using $(\lambda_k-d_i)V_{ik}=(EV)_{ik}$.
At $\lambda=d$ the quotient is $1/(2d)$. With
$u=(\lambda-d)/d$, its stable form is
\begin{equation}
 \frac1d\frac{(1+u)\log(1+u)-u}{u^2}
 =\frac1d\sum_{n=0}^{\infty}\frac{(-u)^n}{(n+1)(n+2)}.
 \label{eq:num-scalar-remainder}
\end{equation}
Ten terms are used for $|u|<0.02$, where the omitted tail is
below double-precision accuracy; otherwise the closed expression
is used. This is a numerical matrix-function identity and
does not impose an interval expansion on the entropy.

For a small interval, take $D$ to be the diagonalized mixture
and $H$ the source density matrix in that basis. Equal traces
give $D^A_{a\mid b}=\sum_i\mathcal R_i$. For the complement,
take the unweighted Gram matrix
\begin{equation}
 H=\begin{pmatrix}\rho_a^A&T_{ab}^A\\
 (T_{ab}^A)^\dagger&\rho_b^A\end{pmatrix}=2\mathcal G^*.
 \label{eq:num-unweighted-gram}
\end{equation}
Diagonalize its two diagonal blocks separately and let $D$
contain their eigenvalues. The perturbation $E$ then has only
off-diagonal blocks. Equation~\eqref{eq:num-gram-entropy} gives
\begin{equation}
 \delta_{a\mid b}=\tr P_a[f(H)-f(D)]
                 =\sum_{i\in a}\mathcal R_i,
 \qquad \delta_{b\mid a}=\sum_{i\in b}\mathcal R_i.
 \label{eq:num-stable-deficit}
\end{equation}
This permits a deficit much smaller than machine precision
relative to $\log2$ to be resolved on its own scale.

The identities are exact on positive support. In the
double-precision implementation, reference eigenvalues below
$\tau=10^{-15}$ are discarded within each symmetry sector.
This is a numerical approximation, checked against
$\tau=10^{-14},10^{-16}$ at sampled endpoint lengths, the
ordinary spectral trace, and the full Gram calculation.
Table~\ref{tab:num-checks} records the supplied production
diagnostics; its cutoff sensitivities refer to retained endpoint
lengths $6,10$ for the dense scan and $10$ for the earlier
extension. Absolute trace-log agreement alone cannot certify
the relative accuracy of an extremely small deficit.
\begin{table}[htbp]\centering\small
\begin{tabular}{lrr}\toprule
Check & Count & Largest error\\\midrule
Unit trace & 6390 & $6.661\times10^{-16}$\\
Hermiticity & 6390 & $1.791\times10^{-77}$\\
Negative eigenvalue magnitude & 6390 & $3.324\times10^{-16}$\\
Small-interval relative cutoff sensitivity & 4824 & $8.635\times10^{-9}$\\
Small-interval trace-log difference (absolute) & 4824 & $2.615\times10^{-13}$\\
Relative local LU residual & 945 & $1.388\times10^{-13}$\\
Gram negative eigenvalue magnitude & 6390 & $1.240\times10^{-16}$\\
Deficit relative cutoff sensitivity & 4824 & $6.852\times10^{-9}$\\
Full-Gram difference (absolute) & 4824 & $1.638\times10^{-12}$\\\bottomrule
\end{tabular}
\caption{Recorded numerical checks on the supplied extended data.
Relative cutoff sensitivity is $\max_\tau|q_\tau-q_{10^{-15}}|/q_{10^{-15}}$.
The LU residual is the Frobenius norm of the residual in
Eq.~\eqref{eq:num-anchor-solve} divided by that of its right-hand side.
Positivity errors are $\max(0,-\lambda_{\min})$. These are numerical
precision checks, not estimates of the lattice-to-CFT error.}
\label{tab:num-checks}
\end{table}

\subsection{Fitting models, residuals, and window dependence}
\label{app:num-fitting}
Write $z=x$, $q=D^A$ for a small interval, and
$z=\ep$, $q=\delta$ for a large interval. For each ordered pair,
we compare
\begin{equation}
 F_{\rm free}(z)=Cz^r,\qquad
 F_1(z)=c_0z^{r_0},\qquad
 F_2(z)=c_0z^{r_0}+c_1z^{r_1}.
 \label{eq:num-fits}
\end{equation}
The free fit estimates $C>0$ and $r$ without theoretical
constraints. $F_1$ and $F_2$ impose the powers in
Table~\ref{tab:num-powers} and fit one and two amplitudes,
respectively; the $F_2$ amplitudes have no sign constraints.
No constant term or extra lattice-correction model is fitted.
The CFT curve uses the analytical leading coefficient, plus
the complete sixth-order term for the small-interval
$I\leftrightarrow\varepsilon$ pair; no parameter is adjusted.
\begin{table}[htbp]\centering\small
\begin{tabular}{llrrrr}\toprule
 & & \multicolumn{2}{c}{Small interval} & \multicolumn{2}{c}{Large interval}\\
CFT & Pair (both directions) & $r_0=\alpha$ & $r_1$ & $r_0=\beta$ & $r_1$\\\midrule
Ising & $I,\sigma$ & 2 & 4 & $1/4$ & $1/2$\\
Ising & $I,\varepsilon$ & 4 & 6 & 2 & 4\\
Ising & $\sigma,\varepsilon$ & 2 & 4 & $1/4$ & $1/2$\\
Boson & $00,ss$ & 2 & 4 & 1 & 2\\
Boson & $00,3s,s$ & 2 & 4 & 5 & 6\\
Boson & $ss,3s,s$ & 2 & 4 & 2 & $3^{*}$\\\bottomrule
\end{tabular}
\caption{Powers used in $F_1,F_2$. The asterisk denotes an
allowed three-point sector whose nonzero total coefficient
has not been established. Except for the two small-interval
Ising sixth-order coefficients, complete next coefficients
are not supplied by the analytical calculation.}
\label{tab:num-powers}
\end{table}

Every fit minimizes unweighted squared residuals in the
observable, $\sum_{i:z_i\in W}[q_i-M(z_i)]^2$, on the window
in Table~\ref{tab:num-windows}. Logarithmic figure axes do not
change this objective. For fixed powers, scale the columns as
$X_{ij}=(z_i/z_\star)^{r_j}$, where $z_\star$ is the upper
cutoff, solve by singular-value decomposition, and convert
the fitted amplitudes back by dividing by $z_\star^{r_j}$.
For a trial free exponent the optimal scaled amplitude is
\begin{equation}
 A(r)=\frac{\sum_i q_i(z_i/z_\star)^r}
                  {\sum_i(z_i/z_\star)^{2r}},\qquad C=A(r)z_\star^{-r}.
 \label{eq:num-profile}
\end{equation}
A scalar search of the profiled objective provides an
independent check on the two-parameter optimization
\cite{numGolubPereyra}. If an amplitude was originally reported
against $(\pi z)^r$, conversion to the present convention uses
$C=\pi^r\widetilde C$ with the fitted exponent, not the
theoretical exponent.

The signed pointwise residual and leading-amplitude discrepancy are
\begin{equation}
 R_i[\%]=100\frac{q_i-M(z_i)}{q_i},\qquad
 \eta_0[\%]=100\left(\frac{c_0}{c_0^{\rm CFT}}-1\right).
 \label{eq:num-residuals}
\end{equation}
Residual ordinates are linear. To test prediction, omit all
observations at one chain size, refit, and predict those
observations. With $M^{(-N_i)}$ denoting the resulting fit,
\begin{align}
 {\rm CVRMSE}_M&=\left\{\frac1m\sum_i
 [q_i-M^{(-N_i)}(z_i)]^2\right\}^{1/2},\nonumber\\
 G_{\rm CV}[\%]&=100\left(1-
          \frac{{\rm CVRMSE}_{F_2}}{{\rm CVRMSE}_{F_1}}\right).
 \label{eq:num-cv}
\end{align}
Positive $G_{\rm CV}$ means improvement. This comparison is
needed because $F_1$ is nested in $F_2$, so training-error
reduction alone is not evidence of better prediction.
The numerical summaries retain the different signs of
coefficient and prediction gains. Parameter fluctuations
under these omissions are sensitivity diagnostics, not
statistical confidence intervals for deterministic data.

The supplied window scan enumerates distinct selections
$0<z\leq v$ through $v=0.20$, requiring at least forty
training observations after every omitted-size refit. For
large intervals there are 1361 admissible Ising selections
and 532 boson selections. The boson endpoint dependence is
reported in Sec.~\ref{sec:num-boson-large}; the adopted
windows are shared across directions rather than chosen
separately for each fit. Finally, $\ell/N\ll1$ and
$\ell\gg1$ are separate continuum requirements. Exact
finite-chain evaluation and numerical precision checks do
not remove finite-$\ell$ errors, and increasing $N$ alone
does not establish convergence of a fitted amplitude to CFT.
